% ============================================================
% SECTION 3: MODELS AND METHODS
% ============================================================
% Target: 1 500--2 000 words.
%
% Cover:
%   3.1  Problem Formulation (self-financing portfolio, loss function)
%   3.2  Asset Price Models  (Black-Scholes; add Heston etc. if applicable)
%   3.3  Neural Network Architecture
%   3.4  Training Procedure
% ============================================================

\section{Models and Methods}
\label{sec:methods}

% ---- 3.1 Problem Formulation ----------------------------------------
\subsection{Problem Formulation}
\label{subsec:problem}

Consider a European-style derivative security with payoff $\Phi(S_T)$ at maturity $T$,
written on a single underlying asset with price process $(S_t)_{t \ge 0}$. A seller of
the derivative charges a premium $\pi$ and seeks a self-financing, discretely-rebalanced
hedging strategy $(\delta_n)_{n=0}^{N-1}$ that minimises the discrepancy between the
terminal portfolio value $V_N$ and the payoff $\Phi(S_N)$.

Trading occurs at dates $0 = t_0 < t_1 < \cdots < t_N = T$. At each time $t_n$ the
hedger holds $\delta_n$ units of the asset; $\delta_n$ must be
$\mathcal{F}_{t_n}$-measurable (no look-ahead). The self-financing condition gives the
discounted portfolio value
\begin{equation}
    \bar{V}_n = \pi + \sum_{j=0}^{n-1} \delta_j \,\Delta\bar{S}_j,
    \qquad \Delta\bar{S}_j := \bar{S}_{j+1} - \bar{S}_j,
    \label{eq:portfolio}
\end{equation}
where $\bar{S}_n = S_n e^{-r t_n}$ is the discounted asset price and $r$ is the
continuously compounded risk-free rate.

The hedging objective is to minimise the loss function
\begin{equation}
    \mathcal{L}(\mathbf{w}) := \frac{1}{M}\sum_{m=1}^{M}
        \Bigl(\Phi\!\left(S_N^{(m)}\right) - V_N^{(m)}\Bigr)^2,
    \label{eq:loss}
\end{equation}
over the neural network weights $\mathbf{w}$, where $M$ is the number of simulated
paths.

% TODO: Mention alternative loss functions if explored (CVaR, entropic risk, etc.)

% ---- 3.2 Asset Price Models -----------------------------------------
\subsection{Asset Price Models}
\label{subsec:models}

\subsubsection{Black--Scholes Model}
\label{subsubsec:bs}

Under the risk-neutral measure $\mathbb{Q}$, the asset price satisfies
\begin{equation}
    \mathrm{d}S_t = r S_t \,\mathrm{d}t + \sigma S_t \,\mathrm{d}W_t,
    \label{eq:bs_sde}
\end{equation}
where $\sigma > 0$ is the constant volatility and $(W_t)$ is a standard Brownian
motion. The exact discretisation used in simulation is
\[
    S_{t_{n+1}} = S_{t_n} \exp\!\Bigl[\Bigl(r - \tfrac{1}{2}\sigma^2\Bigr)\Delta t
        + \sigma\sqrt{\Delta t}\,Z_{n+1}\Bigr], \qquad Z_{n+1}\sim\mathcal{N}(0,1).
\]
The closed-form call price and delta are
\begin{align}
    C_0^{\mathrm{BS}} &= S_0 \mathcal{N}(d_+) - K e^{-rT}\mathcal{N}(d_-),
    \label{eq:bs_price}\\
    \delta_t^{\mathrm{BS}} &= \mathcal{N}(d_+),
    \label{eq:bs_delta}
\end{align}
where $d_{\pm} = \bigl[\ln(S_0/K) + (r \pm \tfrac{1}{2}\sigma^2)T\bigr]/(\sigma\sqrt{T})$
and $\mathcal{N}$ denotes the standard normal CDF. These serve as the theoretical
benchmark against which we compare the neural network hedge.

% TODO: Add further models if applicable, e.g.:
% \subsubsection{Heston Stochastic Volatility Model}
% \label{subsubsec:heston}

% ---- 3.3 Neural Network Architecture --------------------------------
\subsection{Neural Network Architecture}
\label{subsec:architecture}

% TODO: Describe your network architecture in detail.
%   - Type: feedforward / recurrent / weight-sharing.
%   - Inputs at each t_n (e.g. S_n, t_n, moneyness K/S_n, time-to-maturity).
%   - Number of hidden layers, units per layer, activation functions.
%   - How the premium pi is parameterised (trainable scalar / output of a sub-network).
%   - Any constraints on deltas (e.g. clipping, sigmoid output for bounded strategies).
%   - Reference your architecture diagram here if you include one.

Our neural network, inspired by the Deep Hedging framework of \citet{BuehlerEtAl2019},
consists of $N$ sub-networks sharing the same architecture but with independent weights.
At each time step $t_n$, sub-network $\mathcal{N}^{(n)}$ receives the asset price $S_n$
as input and outputs the hedge ratio $\delta_n$. The premium $\pi$ is a separate
trainable scalar initialised at zero.

% TODO: Fill in architectural specifics (layers, units, activations).

% ---- 3.4 Training Procedure -----------------------------------------
\subsection{Training Procedure}
\label{subsec:training}

% TODO: Describe:
%   - Number of paths M, time steps N, batch size, number of epochs.
%   - Optimiser (Adam) and learning rate schedule.
%   - Validation / early stopping / regularisation (dropout, L2, etc.).
%   - Hardware and software (Python version, Keras/TF/PyTorch, Google Colab / local GPU).
%   - How model weights are saved and loaded (checkpointing).

Training was performed by minimising~\eqref{eq:loss} via the Adam optimiser
\citep{KingmaBa2015} with initial learning rate $\eta = \ldots\,$. We generated
$M = 10^5$ paths with $N = 125$ time steps (approximately daily rebalancing over
$T = 0.5$ years). Paths were divided into training ($80\%$), validation ($10\%$),
and test ($10\%$) sets. The network was implemented in Python using Keras
\citep{Chollet2026} with a TensorFlow backend.
